% sec-00: front matter.  Abstract, the no-agreement notice, and the reading path.
% No float.
\section*{Abstract}
\addcontentsline{toc}{section}{Abstract}

A funder buys a trial, and a trial needs a site. This packet accompanies five
letters: two to La Jolla institutions approached in parallel on the same day, two
to disease foundations asked a mechanism question rather than for money, and one
to a clinical trial support service asking for the determination that gates
every other item, which is which party is capable of serving as
sponsor-investigator \cite{ucsdactri2026}.

It sets out the five questions a 45-minute feasibility meeting has to settle, the
obligations that would fall to a site, to the sponsor, and to the agent's
developer, and the site start-up money that is held and not committed.

\textbf{No agreement of any kind exists with UC San Diego, with Moores Cancer
Center, with Scripps Research, with Scripps Health, or with any other
institution.} Both institutions are named at the feasibility stage only. Neither
is described as a partner, sponsor, site, or endorser, and neither is described
to the other as preferred \cite{ucsdmoores2026,scrippsdtc2026}.

\section{What This Packet Is For}

This is the third of five daily packets in
\appfile{funding/auto-fund}. Day 1 re-contacted five federal mechanisms after the
approval of Rasonque on August 26, 2026 \cite{fdarasonque2026}. Day 2 addressed
the private capital the federal route does not reach. This day addresses the
thing both of those exist to fund.

\subsection{The order these three days are in, and why}

A site asks who is funding the work and then asks who the investigator will be.
The first question is answered by day 1 and the second by the conflict boundary
set out in day 2, so approaching a site before either answer existed would have
produced a conversation that stalled on its second question
\cite{ecfr2026part54}.

The same logic governs what is not sent on this day. No engineering detail
travels before a confidentiality agreement is in place, and the route for that
agreement is asked as a question in the fifth letter rather than assumed in the
first.

\subsection{How to read this}

A site principal investigator or a trials office should read \S2, which is the
meeting agenda, and \S3, which is the obligation split. Those two are the whole
of what a first meeting needs.

A foundation reviewer should read \S4, which sets out the funnel and the three
mechanisms, and \S0 above for what is and is not claimed.

A funder or a holder should read \S5, which shows that site start-up money is
allocated, segregated, and committed to nobody. \S1 explains why two institutions
are approached at once, \S6 carries the positioning constraints, the method note,
and the references.
